\chapter*{Вступ}
\addcontentsline{toc}{chapter}{Вступ}

\vspace{-1.2em} % Змінюйте це число, щоб регулювати відступ (чим більше число, тим ближче текст)


Платформи оголошень є одним із найпоширеніших цифрових просторів для купівлі-продажу товарів і послуг, де користувачі щодня вступають у угоди з різними контрагентами. Їхня популярність зумовлена простотою використання, швидким пошуком товарів, низькими бар’єрами входу та широким вибором пропозицій. Однак саме ця відкритість і висока динамічність взаємодії створюють сприятливе середовище для шахрайських практик. На ринку віртуальних угод зловмисники не прагнуть зламати технічні системи, а навпаки — активно використовують людський фактор: бажання придбати товар дешевше, поспіх, довіру до візуальних ознак легітимності та страх втратити вигідну пропозицію.

У структурі шахрайства на платформах оголошень ключовим є не технічний збій, а соціальна інженерія. Зловмисник створює враження реального продавця, покупця, орендодавця або роботодавця, але в реальності використовує підроблені профілі, фальшиві відгуки, привертаючі фото й маніпулятивні формулювання. Така схема працює завдяки тому, що людина не завжди перевіряє контрагента, а приймає рішення під впливом емоцій: зацікавлення, тривоги, відчуття дефіциту або бажання скористатися «вигідною» пропозицією. У результаті навіть досвідчений користувач може стати жертвою, якщо дія відбувається в умовах штучно створеного часу та психологічного тиску.

Сучасні схеми шахрайства на платформах оголошень значно ускладнилися. Якщо раніше домінували прості втрати на передоплаті за неіснуючий товар, то сьогодні в арсеналі зловмисників є комплексні схеми: підроблені доставки, фейкові профілі, імітація перевірки особи, вимагання завдатку за оренду, псевдопослуги, а також шахрайські пропозиції працевлаштування. Часто ці дії комбінуються з фішинговими посиланнями, підробленими вебсайтами та мессенджер-комунікаціями, що дозволяє шахраям не лише викрасти гроші, а й отримати персональні дані, реквізити карток, дані паспорта або доступ до облікових записів. Саме тому аналіз шахрайських схем має не лише практичне, а й теоретичне значення: він дозволяє ідентифікувати типові ознаки загрози, зрозуміти механізми психологічного впливу та сформувати ефективні правила цифрової безпеки.

Наведені далі типи шахрайських схем демонструють, як на практиці виявляються основні принципи обману на онлайн-платформах: створення ілюзії легітимності, емоційне прискорення рішень, використання чужої репутації та формування атмосферу терміновості. Вивчення цих схем має особливу актуальність у контексті розвитку електронної комерції, популяризації маркетплейсів та зростання ролі мобільних додатків у повсякденні. Тому подальший розгляд окремих форм шахрайства є необхідним елементом загальної системи протидії кіберзлочинності та підвищення цифрової грамотності користувачів.


\clearpage
